\documentclass[11pt,a4paper]{article}
\usepackage[utf8]{inputenc}
\usepackage[T1]{fontenc}
\usepackage[french]{babel}
\frenchbsetup{StandardLists=true}
\usepackage[left=2cm,right=2cm,top=2cm,bottom=2cm]{geometry}
\usepackage[dvipsnames]{xcolor}
\usepackage{fouriernc}
\usepackage{amsmath,amssymb,amsfonts}
\usepackage{array,tabularx,booktabs,multirow}
\usepackage{colortbl}
\usepackage{graphicx}
\usepackage{mdframed}
\usepackage{enumitem}
\usepackage{fancyhdr}
\usepackage{chemformula}
\usepackage{awesomebox}
\setlength{\aweboxvskip}{0mm}
\usepackage{tcolorbox}
\tcbuselibrary{skins}
\usepackage{fontawesome5}

\definecolor{exoheader}{HTML}{1E5AA0}
\definecolor{exolight}{HTML}{DCE8FF}
\definecolor{warnorange}{RGB}{220,100,0}
\definecolor{problemebg}{HTML}{FFFDE7}

\renewcommand{\arraystretch}{1.8}
\setlength{\extrarowheight}{1mm}

\pagestyle{fancy}
\renewcommand\headrulewidth{1pt}
\fancyhead[L]{Académie de Besançon}
\fancyhead[C]{\textbf{TP --- Synthèse d'un savon}}
\fancyhead[R]{Terminale / BTS Chimie}
\fancyfoot[C]{page \thepage}
\fancyfoot[R]{2025 -- 2026}
\fancyfoot[L]{Alexis RUIZ}

\newcommand{\sectitle}[1]{%
  \vspace{6pt}
  \noindent\colorbox{exoheader}{\parbox{\dimexpr\linewidth-2\fboxsep}%
    {\color{white}\bfseries\large\strut\quad #1}}
  \vspace{4pt}
}
\newcommand{\subsectitle}[1]{%
  \vspace{4pt}
  \noindent\colorbox{exolight}{\parbox{\dimexpr\linewidth-2\fboxsep}%
    {\color{exoheader}\bfseries\strut\quad #1}}
  \vspace{2pt}
}
\newcounter{question}
\newcommand{\question}{%
  \stepcounter{question}%
  \noindent\textcolor{exoheader}{\textbf{Question \thequestion.}}\quad
}

\begin{document}

\noindent\textbf{NOM :}\hfill\textbf{Prénom :}\hfill\textbf{Classe :}\hfill\phantom{x}\\[0.2cm]
\hrule\vspace{0.2cm}

\begin{center}
  {\large\bfseries\textcolor{exoheader}{Travaux Pratiques de Chimie Organique}}\\[4pt]
  {\LARGE\bfseries Saponification --- Synthèse d'un savon}\\[2pt]
  {\small Académie de Besançon --- Belfort}
\end{center}
\hrule\vspace{6pt}

\begin{tcolorbox}[enhanced, colback=problemebg, colframe=exoheader, boxrule=1.5pt,
  title={\faFlask\quad Problématique},
  fonttitle=\bfseries\color{white},
  attach boxed title to top left={yshift=-2mm, xshift=4mm},
  boxed title style={colback=exoheader, colframe=exoheader}]
\textit{Pendant la Seconde Guerre mondiale, le savon était une denrée rare.
Les particuliers en fabriquaient artisanalement à partir d'huile végétale et de soude.
\textbf{Comment réaliser en laboratoire la synthèse d'un savon ?}}
\end{tcolorbox}

\vspace{6pt}

% ══════════════════════════════════════════════════════════════════════════════
\sectitle{Partie I --- Étude de la réaction}
% ══════════════════════════════════════════════════════════════════════════════

\subsectitle{1. La réaction de saponification}

\begin{mdframed}[backgroundcolor=exolight, linecolor=exoheader, linewidth=1.2pt]
La \textbf{saponification} est la réaction d'un ester (triglycéride d'huile végétale)
avec une base forte (\ch{NaOH}) en milieu aqueux à chaud.
Elle produit du \textbf{savon} (sel d'acide gras) et du \textbf{glycérol}.
La réaction est totale et irréversible.
\end{mdframed}

\vspace{4pt}
\question À l'aide du tableau en annexe, identifier les pictogrammes présents sur les flacons
d'éthanol et de soude. Donner les précautions à prendre. \dotfill

\vspace{0.3cm}\noindent\dotfill

\vspace{4pt}
\question L'équation de saponification de l'oléine (triester) est donnée ci-dessous.
Entourer et nommer les groupes caractéristiques des molécules de réactifs et de produits.

\vspace{0.5cm}
\begin{center}
\begin{mdframed}[backgroundcolor=white, linecolor=exoheader, linewidth=1.5pt]
\[
  \underbrace{(\ch{C17H33COO})_3\ch{C3H5}}_{\text{trioléate de glycérol (oléine)}}
  + 3\,\ch{NaOH}
  \;\longrightarrow\;
  3\,\underbrace{\ch{C17H33COONa}}_{\text{savon (oléate de sodium)}}
  + \underbrace{\ch{C3H8O3}}_{\text{glycérol}}
\]
\end{mdframed}
\end{center}

\vspace{4pt}
\question Pourquoi cette réaction est-elle qualifiée de \og saponification \fg{} ?
Quelle est la différence avec une estérification ? \dotfill

\vspace{0.3cm}\noindent\dotfill

% ══════════════════════════════════════════════════════════════════════════════
\sectitle{Partie II --- Protocole expérimental}
% ══════════════════════════════════════════════════════════════════════════════

\begin{mdframed}[backgroundcolor=orange!15, linecolor=warnorange, linewidth=2pt]
\textbf{\textcolor{warnorange}{\faExclamationTriangle\quad ATTENTION}} ---
La soude concentrée est très corrosive. Porter \textbf{lunettes et gants} obligatoirement.
Ne pas utiliser le savon obtenu pour se laver les mains (caustique).
\end{mdframed}

\vspace{4pt}
\subsectitle{I --- Synthèse du savon (chauffage à reflux)}

\begin{enumerate}[leftmargin=1.5cm, label=\textcolor{exoheader}{\textbf{\arabic*.}}]
  \item Dans un ballon, introduire précisément :
    \begin{itemize}
      \item \textbf{20~mL} d'huile d'arachide (éprouvette graduée),
      \item \textbf{10~mL} d'éthanol à 95\% (éprouvette graduée),
      \item \textbf{20~mL} de solution d'hydroxyde de sodium (éprouvette graduée),
      \item quelques grains de \textbf{pierre ponce}.
    \end{itemize}
  \item Adapter le réfrigérant à boules (ouvrir le robinet d'eau).
  \item Porter à \textbf{ébullition douce pendant 30 minutes}.
  \item Couper le chauffage et laisser refroidir 2 minutes à l'air, puis sous l'eau froide.
\end{enumerate}

\vspace{4pt}
\question Compléter le schéma du montage à reflux ci-dessous.

\vspace{3cm}
\begin{center}
\textit{(espace pour schéma)}
\end{center}
\vspace{1cm}

\subsectitle{II --- Relargage, lavage et filtration}

\begin{enumerate}[leftmargin=1.5cm, label=\textcolor{exoheader}{\textbf{\arabic*.}}]
  \item Verser dans un bécher de 250~mL environ 150~mL d'eau saturée en sel (relargage).
  \item Y ajouter le mélange réactionnel.
  \item Filtrer sur Büchner.
  \item Laver le solide avec de l'eau distillée, puis mettre à sécher sur une coupelle.
\end{enumerate}

\vspace{4pt}
\question Quel est le rôle de l'eau saturée en sel (relargage) ? \dotfill

\vspace{0.3cm}\noindent\dotfill

% ══════════════════════════════════════════════════════════════════════════════
\newpage
\sectitle{Partie III --- Propriétés du savon obtenu}
% ══════════════════════════════════════════════════════════════════════════════

\subsectitle{1. Tests sur le savon}

\question Observer l'aspect du savon obtenu. Le décrire. \dotfill

\vspace{0.3cm}\noindent\dotfill

\question Dissoudre quelques grammes de savon dans de l'eau distillée.
Mesurer le pH de la solution. Expliquer pourquoi la solution est basique. \dotfill

\vspace{0.3cm}\noindent\dotfill

\vspace{0.3cm}\noindent\dotfill

\question Tester l'action émulsifiante du savon : ajouter quelques gouttes d'huile dans
un tube à essai contenant de l'eau + savon, et agiter. Que constate-t-on ? \dotfill

\vspace{0.3cm}\noindent\dotfill

\subsectitle{2. Structure moléculaire et pouvoir détergent}

\question La molécule de savon (oléate de sodium) possède une partie hydrophile et
une partie hydrophobe. Identifier ces deux parties sur la formule :

\vspace{0.3cm}
\begin{center}
\ch{C17H33COO^{-}} \ch{Na^{+}}
\end{center}

\vspace{2cm}

\question Expliquer, à l'aide du schéma d'une micelle, comment le savon permet
de dégraisser. \dotfill

\vspace{0.3cm}\noindent\dotfill

\vspace{0.3cm}\noindent\dotfill

% ══════════════════════════════════════════════════════════════════════════════
\newpage
\sectitle{Annexe --- Données}
% ══════════════════════════════════════════════════════════════════════════════

\vspace{4pt}
\begin{center}
\begin{tabularx}{\linewidth}{|>{\centering\bfseries\color{white}\columncolor{exoheader}}m{2.5cm}
                              |>{\centering\arraybackslash}X
                              |>{\centering\arraybackslash}m{3cm}
                              |>{\centering\arraybackslash}m{3cm}|}
\hline
\rowcolor{exoheader}
\color{white}\bfseries Espèce &
\color{white}\bfseries Formule &
\color{white}\bfseries Masse molaire (g/mol) &
\color{white}\bfseries Densité \\
\hline
Oléine (huile) & \ch{(C17H33COO)3C3H5} & 884 & 0,91 \\
\hline
Éthanol & \ch{C2H5OH} & 46 & 0,79 \\
\hline
Soude & \ch{NaOH} & 40 & --- \\
\hline
Oléate de sodium & \ch{C17H33COONa} & 304 & --- \\
\hline
Glycérol & \ch{C3H8O3} & 92 & 1,26 \\
\hline
\end{tabularx}
\end{center}

\end{document}
